\subsection{Введение}

Программирование световой индикации — это не просто способ украсить полёт, но и важнейший инструмент отладки и взаимодействия с пилотом. Светодиоды позволяют дрону «общаться»: сообщать о текущем режиме работы, уровне заряда батареи, готовности к взлёту или возникновении ошибок.

В этом уроке мы пройдём путь от простого включения лампочки до создания сложных динамических световых эффектов, работающих параллельно с полётной программой.

\begin{infobox}[Цели урока]
Мы научимся:
\begin{itemize}
    \item Управлять цветом и яркостью светодиодов через код на Lua.
    \item Использовать таймеры для создания анимаций без задержки выполнения программы.
    \item Проектировать алгоритмы индикации, реагирующие на события дрона.
\end{itemize}
\end{infobox}
